\subsection{NON-SPLIT CASE}

PROPOSITION 14. Let \(k'\) be an extension of \(k\) and \(\mathfrak{g}' = \mathfrak{g} \otimes_k k'\) . Let \(\mathfrak{m}\) be a subalgebra of \(\mathfrak{g}\) and \(\mathfrak{m}' = \mathfrak{m} \otimes_k k'\) . If \(\mathfrak{m}'\) is a parabolic (resp. Borel) subalgebra of \(\mathfrak{g}'\) , \(\mathfrak{m}\) is a parabolic (resp. Borel) subalgebra of \(\mathfrak{g}\) .

By Prop. 8 and 11, it suffices to prove that m contains a splitting Cartan subalgebra of g. Let b be a Borel subalgebra of g. Then \(b' = b \otimes_{k} k'\) is a Borel subalgebra of \(\mathfrak{g}'\) , so \(\mathfrak{m}' \cap \mathfrak{b}'\) contains a Cartan subalgebra of \(\mathfrak{g}'\) (Prop. 10). Let \(\mathfrak{t}\) be a Cartan subalgebra of \(\mathfrak{m} \cap \mathfrak{b}\) . Then \(\mathfrak{t} \otimes_k k'\) is a Cartan subalgebra of \(\mathfrak{m}' \cap \mathfrak{b}'\) , and hence of \(\mathfrak{g}'\) (Chap. VII, §3, no. 3, Prop. 3). Consequently, \(\mathfrak{t}\) is a Cartan subalgebra of \(\mathfrak{g}\) , and it is splitting since it is contained in \(\mathfrak{b}\) .

DEFINITION 3. Let \(\mathfrak{a}\) be a semi-simple (or more generally reductive) Lie algebra and \(\bar{k}\) an algebraic closure of \(k\) . A subalgebra \(\mathfrak{m}\) of \(\mathfrak{a}\) is said to be parabolic (resp. Borel) if \(\mathfrak{m} \otimes_k \bar{k}\) is a parabolic (resp. Borel) subalgebra of \(\mathfrak{a} \otimes_k \bar{k}\) .

If \(\mathfrak{a}\) is splittable, Prop. 14 shows that this definition is equivalent to Definition 2 (resp. to Definition 1).

PROPOSITION 15. Let \(\mathfrak{a}\) be a reductive Lie algebra, \(k'\) an extension of \(k\) , and \(\mathfrak{m}\) a subalgebra of \(\mathfrak{a}\) . Then \(\mathfrak{m}\) is a parabolic (resp. Borel) subalgebra of \(\mathfrak{a}\) if and only if \(\mathfrak{m} \otimes_k k'\) is a parabolic (resp. Borel) subalgebra of \(\mathfrak{a} \otimes_k k'\) .

This follows immediately from Prop. 14.

